\section{Podsumowanie możliwości silnika}

SabberStone jest wszechstronnym środowiskiem badawczym, które pokrywa potrzeby eksperymentalne na wielu poziomach --- od szybkiego prototypowania algorytmów po organizację formalnych turniejów i integrację z zewnętrznymi systemami uczenia maszynowego. Tabela~\ref{tab:sabber_summary} zestawia kluczowe możliwości silnika.

\begin{table}[ht]
\centering
\caption{Zestawienie kluczowych możliwości silnika SabberStone}
\label{tab:sabber_summary}
\begin{tabular}{p{5.2cm}p{8.3cm}}
\toprule
\textbf{Możliwość} & \textbf{Opis} \\
\midrule
Wysokie pokrycie kart
  & Ponad 98\% kart formatu standardowego (rok 2019), z testami jednostkowymi dla każdej implementacji. \\[4pt]
Częściowa obserwowalność
  & Karty przeciwnika na ręce i w talii są ukryte, co odpowiada warunkom rzeczywistej rozgrywki. \\[4pt]
Model przejść
  & Symulacja przyszłego stanu gry bez modyfikowania bieżącego stanu, co jest podstawą algorytmów przeszukiwania. \\[4pt]
Klonowanie stanu gry
  & Tworzenie niezależnej kopii stanu gry umożliwiającej równoległe eksplorowanie gałęzi drzewa stanów. \\[4pt]
Lista legalnych akcji
  & Silnik automatycznie udostępniania wszystkie dopuszczalne ruchy z pełnym uwzględnieniem reguł gry. \\[4pt]
Interfejs agenta
  & Abstrakcyjny, pięcioetapowy cykl życia agenta z wyraźnym podziałem na inicjalizację, decyzję i finalizację. \\[4pt]
Organizacja turniejów
  & Format round-robin z konfigurowalną liczbą gier na parę agentów i automatyczną agregacją wyników. \\[4pt]
Dwie ścieżki eksperymentalne
  & Gra gotowymi taliami lub samodzielny dobór talii jako część problemu optymalizacji. \\[4pt]
Integracja gRPC
  & Możliwość implementacji agentów w dowolnym języku programowania, którzy mogą komunikować się z silnikiem przez sieć. \\[4pt]
Klient Unity3D
  & Graficzna wizualizacja rozgrywki w czasie rzeczywistym. \\[4pt]

\bottomrule
\end{tabular}
\end{table}
